% ============================================================
% BAB 2 — KAJIAN PUSTAKA
% Scope: Backend / CNN MobileNet / TensorFlow.js / Data Logging / K-Means / Dashboard
% Proyek: Escape the Sketchbook — Interactive HITL AI Literacy Simulation
% Penulis: Muhammad Dias Al-Izzat (Dias)
% ============================================================

\chapter{KAJIAN PUSTAKA}
\label{chap:bab2}

Bab ini menyajikan kajian pustaka yang menjadi landasan teoretis bagi perancangan komponen \textit{backend}, klasifikasi, pencatatan data, dan analisis pada sistem \textit{Escape the Sketchbook}. Kajian dimulai dari deskripsi permasalahan yang menjelaskan mengapa klasifikasi berlabel tunggal belum memadai secara pedagogis, mengapa \textit{confidence score} harus divisualisasikan, dan mengapa pencatatan interaksi harus terstruktur untuk keperluan analisis. Dilanjutkan dengan teori penunjang yang mendasari setiap keputusan teknis—mulai dari literasi kecerdasan buatan, \textit{Human-in-the-Loop}, klasifikasi sketsa, arsitektur CNN dan MobileNet, inferensi \textit{client-side} dengan TensorFlow.js, \textit{confidence score}, \textit{data logging}, hingga \textit{K-Means clustering}—serta ditutup dengan penelitian terkait dan analisis \textit{state of the art} yang memosisikan kontribusi teknis proyek ini di antara kajian-kajian sebelumnya.

% ============================================================
\section{DESKRIPSI PERMASALAHAN}
\label{sec:deskripsi-permasalahan}
% ============================================================

Permasalahan utama yang menjadi landasan proyek ini adalah belum adanya media interaktif yang memfasilitasi siswa SMP kelas 7--9 untuk mengevaluasi output kecerdasan buatan secara reflektif melalui simulasi klasifikasi sketsa. Sebagian besar aplikasi klasifikasi gambar yang tersedia saat ini menampilkan satu label prediksi sebagai output final, tanpa menunjukkan distribusi probabilitas di balik keputusan tersebut. Pendekatan klasifikasi berlabel tunggal ini secara pedagogis tidak memadai karena menyembunyikan ketidakpastian inheren dari model—siswa tidak mendapat kesempatan untuk mengamati bahwa AI seringkali tidak ``yakin'' dengan satu jawaban, melainkan mendistribusikan probabilitas ke beberapa kelas sekaligus [1]. Tanpa visibilitas terhadap distribusi probabilitas tersebut, siswa cenderung menerima output AI secara \textit{wholesale}, yang merupakan manifestasi dari fenomena \textit{automation bias}—kecenderungan manusia untuk mempercayai output sistem otomatis tanpa evaluasi kritis [4].

Persoalan ini diperkuat oleh temuan dalam literatur \textit{Explainable AI} (XAI) di konteks pendidikan. Khosravi \textit{et al.} [3] menekankan bahwa transparansi proses keputusan AI merupakan prasyarat bagi peserta didik untuk dapat menginterrogasi dan mengevaluasi output tersebut. Dalam konteks klasifikasi sketsa, transparansi ini dapat diwujudkan melalui visualisasi \textit{confidence score} pada \textit{Top-3 prediction}, sehingga siswa dapat melihat bukan hanya apa yang ``dipilih'' oleh AI, tetapi juga seberapa yakin AI terhadap pilihannya dan alternatif-alternatif apa yang dipertimbangkan. Visualisasi \textit{confidence score} bukan sekadar fitur teknis, melainkan kebutuhan pedagogis yang memungkinkan siswa memiliki data untuk mengevaluasi output AI secara lebih bermakna.

Selain permasalahan visibilitas output, terdapat pula tantangan dalam hal pencatatan interaksi siswa. Agar pola keputusan siswa terhadap output AI dapat dianalisis secara sistematis, diperlukan mekanisme \textit{data logging} yang terstruktur—bukan sekadar mencatat hasil akhir, melainkan menangkap konteks lengkap setiap keputusan: sketsa yang diklasifikasikan, distribusi \textit{confidence score} dari model, keputusan yang diambil siswa (setuju atau menolak prediksi AI), serta waktu respons. Pencatatan terstruktur ini esensial untuk analisis \textit{post-hoc} menggunakan metode seperti \textit{K-Means clustering}, yang memerlukan data numerik yang konsisten dan terdefinisi untuk menghasilkan kelompok pola keputusan yang dapat diinterpretasikan [5][10]. Tanpa \textit{logging} terstruktur, data interaksi tidak dapat digunakan untuk mengidentifikasi pola—misalnya, apakah seorang siswa cenderung menerima prediksi AI tanpa mempertimbangkan alternatif, atau justru mengevaluasi \textit{confidence score} sebelum mengambil keputusan.

Dengan demikian, permasalahan ini memiliki tiga dimensi yang saling terkait: (1) kebutuhan akan klasifikasi multi-label dengan visibilitas \textit{confidence} untuk memerangi \textit{automation bias}, (2) kebutuhan akan \textit{logging} terstruktur sebagai fondasi analisis perilaku keputusan, dan (3) kebutuhan akan metode analisis yang mampu mengelompokkan pola keputusan siswa secara interpretif. Proyek \textit{Escape the Sketchbook} merespons ketiga dimensi ini melalui integrasi klasifikasi sketsa berbasis MobileNet dengan visualisasi \textit{Top-3 confidence}, mekanisme \textit{Human-in-the-Loop} yang mencatat keputusan siswa, serta analisis \textit{K-Means clustering} pada data log untuk mengidentifikasi pola keputusan.

% ============================================================
\section{TEORI PENUNJANG}
\label{sec:teori-penunjang}
% ============================================================

Bagian ini membahas teori-teori yang mendasari perancangan komponen teknis pada proyek \textit{Escape the Sketchbook}, meliputi literasi kecerdasan buatan, \textit{Human-in-the-Loop} pada sistem klasifikasi, klasifikasi sketsa, arsitektur CNN dan MobileNet, TensorFlow.js dan inferensi \textit{client-side}, \textit{confidence score}, \textit{data logging}, serta \textit{K-Means clustering}.

% ============================================================
\subsection{Literasi Kecerdasan Buatan}
\label{subsec:literasi-kb}
% ============================================================

Literasi kecerdasan buatan (\textit{AI literacy}) didefinisikan oleh Ng \textit{et al.} [1] sebagai seperangkat kompetensi yang memungkinkan individu untuk memahami, menggunakan, mengevaluasi, dan merefleksikan teknologi AI dalam konteks kehidupan sehari-hari. Dalam kerangka konseptual yang mereka ajukan, terdapat empat dimensi utama: \textit{know and understand AI}, \textit{use and apply AI}, \textit{evaluate and create AI}, serta \textit{AI ethics}. Dimensi evaluasi menjadi sangat relevan dalam konteks proyek ini, karena mencakup kemampuan untuk menilai output AI secara kritis—termasuk memahami bahwa output AI bersifat probabilistik, bukan deterministik. Ng \textit{et al.} secara eksplisit menekankan bahwa penalaran probabilistik (\textit{probabilistic reasoning}) merupakan komponen fundamental dari literasi AI: individu yang melek AI harus mampu mengevaluasi output probabilistik dari model, memahami bahwa \textit{confidence score} merepresentasikan tingkat keyakinan model—bukan kebenaran absolut—dan mempertimbangkan alternatif-alternatif yang mungkin sama masuk akalnya [1].

Dalam konteks K-12, Casal-Otero \textit{et al.} [2] melakukan tinjauan sistematis terhadap implementasi literasi AI di jenjang pendidikan dasar dan menengah. Mereka menemukan bahwa strategi literasi AI yang efektif untuk siswa usia sekolah harus mempertimbangkan tahap perkembangan kognitif peserta didik—siswa SMP belum sepenuhnya mampu berpikir abstrak tentang konsep probabilistik, sehingga diperlukan strategi yang sesuai usia (\textit{age-appropriate strategies}) untuk membuat konsep tersebut teramati dan bermakna. Salah satu pendekatan yang direkomendasikan adalah menggunakan contoh konkret di mana siswa dapat melihat distribusi probabilitas secara langsung, alih-alih menjelaskan konsep probabilistik secara abstrak [2].

Koneksi antara teori literasi AI dan proyek ini terletak pada mekanisme klasifikasi sketsa dengan visualisasi \textit{Top-3 confidence score}. Melalui mekanisme ini, konsep probabilistik yang abstrak menjadi teramati: ketika model MobileNet mengklasifikasikan sketsa siswa dan menampilkan tiga prediksi teratas beserta nilai \textit{confidence}-nya, siswa dapat secara langsung mengamati bahwa AI mendistribusikan keyakinannya ke beberapa kelas—bukan hanya menunjuk satu jawaban dengan kepastian penuh. Dengan demikian, klasifikasi sketsa dalam \textit{Escape the Sketchbook} berfungsi sebagai contoh konkret yang membuat penalaran probabilistik terlihat dan dapat dievaluasi oleh siswa SMP, sejalan dengan rekomendasi Casal-Otero \textit{et al.} [2] untuk strategi literasi AI yang sesuai usia.

% ============================================================
\subsection{Human-in-the-Loop pada Sistem Klasifikasi}
\label{subsec:hitl-klasifikasi}
% ============================================================

Konsep \textit{Human-in-the-Loop} (HITL) merujuk pada paradigma di mana manusia berpartisipasi secara aktif dalam alur keputusan sistem otomatis, baik dengan memberikan masukan, mengoreksi output, maupun membuat keputusan final [4]. Mosqueira-Rey \textit{et al.} [4] dalam tinjauan komprehensif mereka mengidentifikasi bahwa HITL menjadi penting terutama ketika output sistem otomatis memiliki konsekuensi signifikan dan ketika terdapat risiko \textit{automation bias}—yaitu kecenderungan manusia untuk menerima output sistem otomatis secara tidak kritis karena asumsi bahwa mesin ``lebih tahu''. Dalam konteks klasifikasi, \textit{automation bias} terjadi ketika pengguna menerima label prediksi AI tanpa mempertimbangkan kemungkinan kesalahan atau alternatif lain, yang justru menghilangkan kesempatan untuk evaluasi kritis [4].

Memarian \& Doleck [5] melakukan analisis khusus terhadap peran HITL dalam domain \textit{Artificial Intelligence in Education} (AIED). Mereka menemukan bahwa dalam sistem AIED, peran manusia dapat dikategorikan menjadi beberapa tipe: sebagai \textit{supervisor} yang mengawasi keputusan sistem, sebagai \textit{corrector} yang memperbaiki output yang salah, dan sebagai \textit{decision-maker} yang membuat keputusan final berdasarkan informasi dari sistem [5]. Temuan penting dari kajian mereka adalah bahwa efektivitas HITL bergantung pada kualitas informasi yang diberikan kepada manusia—jika sistem hanya menampilkan output final tanpa konteks, manusia cenderung menjadi \textit{rubber stamp} yang sekadar menyetujui keputusan mesin. Sebaliknya, ketika sistem menyediakan informasi yang memungkinkan evaluasi—seperti distribusi \textit{confidence score}—manusia memiliki data yang diperlukan untuk membuat penilaian yang lebih bermakna [5].

Dalam konteks proyek \textit{Escape the Sketchbook}, implementasi HITL dirancang dengan prinsip yang jelas: sistem menyediakan informasi klasifikasi (\textit{Top-3 confidence score}) kepada siswa, kemudian siswa membuat keputusan apakah setuju atau menolak prediksi AI tersebut. Keputusan siswa dicatat sebagai \textit{log} untuk keperluan analisis, namun \textbf{tidak digunakan untuk melatih ulang model}. Pernyataan ini penting untuk membedakan paradigma HITL yang diterapkan dalam proyek ini dari paradigma \textit{active learning} yang menggunakan umpan balik manusia untuk memperbarui parameter model. Dalam proyek ini, model MobileNet bersifat statis—dilatih sebelumnya dengan \textit{dataset} sketsa dan dijalankan di \textit{client-side} tanpa pembaruan bobot—sementara keputusan siswa dicatat murni sebagai data observasi untuk analisis pola perilaku. Dengan demikian, HITL dalam sistem ini berfungsi sebagai mekanisme pencatatan keputusan (\textit{decision logging}), bukan mekanisme pelatihan ulang (\textit{model retraining}).

% ============================================================
\subsection{Klasifikasi Sketsa}
\label{subsec:klasifikasi-sketsa}
% ============================================================

Klasifikasi sketsa (\textit{sketch classification}) merupakan tugas mengkategorikan gambar sketsa tangan ke dalam kelas-kelas yang telah ditentukan. Meskipun secara konseptual serupa dengan klasifikasi gambar fotografis, klasifikasi sketsa memiliki karakteristik tersendiri yang menjadikannya lebih menantang. Xu \textit{et al.} [6] dalam survei komprehensif mereka mengenai \textit{deep learning} untuk sketsa tangan bebas (\textit{free-hand sketch}) mengidentifikasi beberapa perbedaan fundamental antara sketsa dan foto: (1) sketsa bersifat \textit{sparse}—hanya berisi garis-garis esensial tanpa informasi warna atau tekstur, (2) sketsa memiliki tingkat abstraksi yang bervariasi—dua orang dapat menggambar objek yang sama dengan level detail yang sangat berbeda, dan (3) sketsa seringkali ambigu—sebuah sketsa yang dimaksudkan sebagai ``bintang'' dapat menyerupai ``roda'' atau ``matahari'' tergantung pada interpretasi penggambar [6].

Tantangan klasifikasi sketsa semakin kompleks dalam konteks siswa SMP karena variasi gaya menggambar yang dipengaruhi oleh faktor usia dan kemampuan motorik. Sketsa yang dibuat oleh siswa kelas 7 cenderung berbeda dalam tingkat detail dan proporsi dibandingkan sketsa yang dibuat oleh siswa kelas 9, meskipun keduanya menggambar objek yang sama. Variasi ini menyebabkan distribusi \textit{confidence score} pada model klasifikasi menjadi lebih menyebar—model seringkali tidak dapat memberikan prediksi dengan \textit{confidence} tinggi pada satu kelas saja, melainkan mendistribusikan probabilitas ke beberapa kelas sekaligus. Kondisi ini justru menguntungkan secara pedagogis, karena siswa dapat melihat bahwa AI tidak selalu ``yakin'' dan perlu mengevaluasi alternatif-alternatif yang ditampilkan.

Dalam proyek \textit{Escape the Sketchbook}, klasifikasi sketsa tidak hanya berfungsi sebagai tugas pengenalan pola, tetapi juga sebagai sarana untuk memperlihatkan proses evaluasi terhadap output AI. Kategori sketsa yang digunakan dalam sistem ini terbagi menjadi tiga subset—\textit{Solid} (benda-benda padat seperti meja, kursi, batu), \textit{Danger} (benda-benda berbahaya seperti api, pisau, gunting), dan \textit{Decorative} (benda-benda dekoratif seperti bunga, kupu-kupu, pelangi)—yang masing-masing memiliki konsekuensi berbeda dalam simulasi. Pembagian ke dalam subset kategori ini diperlukan karena klasifikasi multi-kelas pada level benda individual (misalnya 30+ kelas) menghasilkan \textit{confidence score} yang sangat tersebar dan sulit diinterpretasikan oleh siswa SMP, sementara pengelompokan ke dalam tiga kategori yang lebih kasar (\textit{coarse-grained}) menghasilkan distribusi \textit{confidence} yang lebih terbaca dan bermakna untuk keperluan evaluasi oleh siswa.

% ============================================================
\subsection{Convolutional Neural Network dan MobileNet}
\label{subsec:cnn-mobilenet}
% ============================================================

\textit{Convolutional Neural Network} (CNN) merupakan arsitektur jaringan saraf tiruan yang dirancang khusus untuk pemrosesan data berbasis grid, seperti gambar. Arsitektur CNN memanfaatkan operasi konvolusi untuk mengekstraksi fitur hierarkis dari input: lapisan-lapisan awal menangkap fitur-fitur dasar seperti tepi dan tekstur, sementara lapisan-lapisan lebih dalam menangkap pola-pola kompleks seperti bentuk dan struktur objek. Kemampuan CNN untuk mempelajari representasi fitur secara otomatis dari data—tanpa memerlukan \textit{feature engineering} manual—menjadikannya arsitektur dominan untuk tugas-tugas klasifikasi gambar, termasuk klasifikasi sketsa [6].

MobileNet merupakan varian CNN yang dikembangkan oleh Google dengan tujuan utama menghasilkan model yang ringan (\textit{lightweight}) dan efisien untuk perangkat dengan sumber daya terbatas. Efisiensi MobileNet dicapai melalui dua mekanisme utama: (1) \textit{depthwise separable convolution}, yang memecah operasi konvolusi standar menjadi \textit{depthwise convolution} dan \textit{pointwise convolution}, sehingga secara signifikan mengurangi jumlah parameter dan operasi komputasi, serta (2) dua \textit{hyperparameter}—\textit{width multiplier} ($\alpha$) dan \textit{resolution multiplier} ($\rho$)—yang memungkinkan penyesuaian \textit{trade-off} antara akurasi dan efisiensi. Dengan \textit{width multiplier} $\alpha < 1$, jumlah \textit{channel} pada setiap lapisan dikurangi secara proporsional, menghasilkan model yang lebih kecil dan lebih cepat dengan pengurangan akurasi yang relatif minimal [6].

Penggunaan MobileNet dalam proyek ini didasarkan pada pertimbangan bahwa model harus dijalankan di \textit{client-side}—yaitu di \textit{browser} siswa melalui TensorFlow.js—tanpa bergantung pada server inferensi. Karena inferensi berjalan di perangkat siswa yang bervariasi (mulai dari laptop spesifikasi rendah hingga perangkat yang lebih bertenaga), model yang digunakan harus memiliki ukuran file yang cukup kecil untuk diunduh ke \textit{browser} dan kebutuhan komputasi yang rendah cukup untuk menghasilkan prediksi dalam waktu yang dapat ditoleransi. MobileNet, dengan ukuran model sekitar beberapa megabyte dan kebutuhan komputasi yang jauh lebih rendah dibandingkan arsitektur seperti ResNet atau VGG, memenuhi persyaratan ini.

Selain itu, pendekatan \textit{transfer learning} diterapkan untuk mengadaptasi model MobileNet yang telah dilatih pada \textit{dataset} ImageNet ke tugas klasifikasi sketsa. \textit{Transfer learning} memungkinkan pemanfaatan fitur-fitur visual yang telah dipelajari dari \textit{dataset} sumber (misalnya, deteksi tepi, tekstur, dan bentuk dasar) untuk tugas target yang memiliki jumlah data latih terbatas. Dalam konteks ini, lapisan-lapisan \textit{base} MobileNet yang telah dilatih sebelumnya digunakan sebagai \textit{feature extractor}, sementara lapisan-lapisan \textit{classifier} di bagian atas diganti dan dilatih ulang menggunakan \textit{dataset} sketsa. Pendekatan ini mengurangi kebutuhan data latih dan waktu pelatihan secara signifikan, sekaligus mempertahankan kualitas representasi fitur yang telah dipelajari.

% ============================================================
\subsection{TensorFlow.js dan Client-Side Inference}
\label{subsec:tfjs-client-side}
% ============================================================

TensorFlow.js merupakan \textit{library} \textit{open-source} yang memungkinkan pelatihan dan inferensi model \textit{machine learning} langsung di \textit{browser} web maupun \textit{Node.js}. Smilkov \textit{et al.} [8] menjelaskan bahwa TensorFlow.js dirancang untuk membawa kemampuan \textit{machine learning} ke ekosistem web tanpa memerlukan instalasi perangkat lunak tambahan atau akses ke server komputasi. Arsitektur TensorFlow.js mendukung dua \textit{backend} utama: \textit{WebGL backend} yang memanfaatkan GPU melalui API \textit{WebGL} untuk akselerasi komputasi, dan \textit{CPU backend} yang berjalan di \textit{CPU} murni sebagai \textit{fallback} ketika WebGL tidak tersedia. Dukungan terhadap WebGL memungkinkan inferensi model yang cukup cepat di \textit{browser}, meskipun tidak secepat inferensi \textit{native} di \textit{GPU} melalui CUDA [8].

Goh, Ho \& Abas [7] melakukan studi komprehensif mengenai penerapan \textit{deep learning} pada \textit{front-end web application} dan mengidentifikasi beberapa pertimbangan kritis dalam proses \textit{deployment}: (1) ukuran model—model yang terlalu besar akan memerlukan waktu unduh yang lama dan mengonsumsi memori \textit{browser} yang signifikan, (2) latensi inferensi—waktu yang dibutuhkan model untuk menghasilkan prediksi harus dalam batas yang dapat ditoleransi untuk pengalaman pengguna yang responsif, dan (3) kompatibilitas \textit{browser}—tidak semua \textit{browser} mendukung WebGL dengan tingkat performa yang sama, sehingga diperlukan strategi \textit{fallback} [7]. Mereka juga menemukan bahwa arsitektur model ringan seperti MobileNet merupakan pilihan yang paling sesuai untuk \textit{deployment front-end}, karena menawarkan keseimbangan terbaik antara akurasi dan efisiensi di lingkungan \textit{browser} [7].

Dalam proyek \textit{Escape the Sketchbook}, inferensi dilaksanakan sepenuhnya di \textit{client-side}—yaitu di \textit{browser} siswa—menggunakan TensorFlow.js. Keputusan arsitektural ini didasarkan pada alasan bahwa sistem tidak memiliki server AI, tidak menggunakan \textit{ML server}, dan tidak bergantung pada \textit{cloud inference}. Server dalam arsitektur proyek ini hanya berperan sebagai \textit{REST API} untuk pencatatan \textit{log} interaksi ke \textit{database} SQLite—bukan untuk menjalankan inferensi model. Dengan demikian, seluruh proses klasifikasi sketsa, mulai dari penerimaan input gambar, pra-pemrosesan, inferensi model MobileNet, hingga penghitungan \textit{confidence score}, berjalan di \textit{browser} siswa. Arsitektur ini memiliki implikasi: di satu sisi, eliminasi ketergantungan pada server inferensi mengurangi latensi jaringan dan masalah privasi data (sketsa siswa tidak perlu dikirim ke server); di sisi lain, performa inferensi bergantung pada kemampuan perangkat siswa, sehingga pemilihan model ringan seperti MobileNet menjadi keharusan arsitektural, bukan sekadar preferensi.

\begin{figure}[H]
\centering
\includegraphics[width=\textwidth]{placeholder}
\caption{Arsitektur inferensi \textit{client-side}: model MobileNet berjalan di \textit{browser} siswa melalui TensorFlow.js, sementara server hanya menangani \textit{REST API logging} ke SQLite.}
\label{fig:arsitektur-client-side}
\end{figure}

% ============================================================
\subsection{Confidence Score}
\label{subsec:confidence-score}
% ============================================================

\textit{Confidence score} merupakan nilai numerik yang merepresentasikan tingkat keyakinan model terhadap setiap kelas prediksi. Dalam konteks klasifikasi menggunakan \textit{softmax output layer}, \textit{confidence score} dihasilkan melalui fungsi \textit{softmax} yang mengkonversi \textit{logit} mentah menjadi distribusi probabilitas di mana jumlah seluruh nilai \textit{confidence} sama dengan satu. Kelas dengan \textit{confidence score} tertinggi menjadi prediksi utama model, namun nilai-nilai \textit{confidence} pada kelas-kelas lain juga mengandung informasi penting—mereka menunjukkan seberapa besar model mempertimbangkan alternatif-alternatif tersebut sebagai kemungkinan yang masuk akal.

Pemahaman tentang \textit{confidence score} memerlukan pembedaan yang krusial antara \textit{confidence} dan \textit{correctness}. \textit{Confidence score} yang tinggi tidak menjamin bahwa prediksi model benar—model dapat sangat yakin namun salah (\textit{high confidence, low correctness}), sebaliknya model dapat kurang yakin namun prediksinya benar (\textit{low confidence, high correctness}). Kesenjangan antara \textit{confidence} dan \textit{correctness} ini merupakan fenomena yang terdokumentasi dengan baik dalam literatur dan merupakan salah satu alasan mengapa visualisasi \textit{confidence} saja tidak cukup—pengguna juga memerlukan konteks untuk menginterpretasikan nilai \textit{confidence} tersebut [9].

Karran, Demazure \& Hudelot [9] dalam kajian mereka mengenai desain visualisasi \textit{confidence} menemukan bahwa representasi \textit{confidence} yang efektif membantu pengguna menimbang output AI secara lebih bermakna. Mereka mengidentifikasi bahwa visualisasi \textit{confidence} yang hanya menampilkan satu nilai (misalnya, ``AI 95\% yakin ini kucing'') cenderung membuat pengguna fokus pada angka tersebut secara terisolasi, sementara visualisasi yang menampilkan distribusi \textit{confidence} di beberapa kelas—seperti \textit{bar chart} untuk \textit{Top-3 prediksi}—memungkinkan pengguna untuk membandingkan dan mengevaluasi alternatif secara relatif [9].

Dalam proyek ini, \textit{confidence score} digunakan dalam dua kapasitas: (1) sebagai informasi yang divisualisasikan kepada siswa melalui \textit{Top-3 prediction}, sehingga siswa dapat melihat distribusi keyakinan model dan mengevaluasi alternatif-alternatif yang dipertimbangkan AI, serta (2) sebagai fitur numerik yang dicatat dalam \textit{log} interaksi untuk keperluan analisis. Salah satu metrik turunan yang dapat diekstraksi dari \textit{confidence score} adalah \textit{confidence gap}—selisih antara \textit{confidence score} tertinggi dan \textit{confidence score} kedua tertinggi. \textit{Confidence gap} yang kecil menunjukkan bahwa model ragu-ragu antara dua kelas, sementara \textit{confidence gap} yang besar menunjukkan bahwa model cukup yakin dengan prediksi utamanya. Metrik ini relevan untuk analisis perilaku siswa: bagaimana siswa merespons situasi ketika model ragu (\textit{confidence gap} kecil) versus ketika model yakin (\textit{confidence gap} besar) memberikan indikasi tentang seberapa kritis siswa mengevaluasi output AI.

% ============================================================
\subsection{Data Logging}
\label{subsec:data-logging}
% ============================================================

\textit{Data logging} dalam konteks sistem interaktif merujuk pada pencatatan terstruktur atas peristiwa-peristiwa (\textit{events}) yang terjadi selama interaksi pengguna dengan sistem. Pencatatan ini mencakup tidak hanya hasil akhir interaksi, tetapi juga konteks, urutan, dan waktu terjadinya setiap peristiwa. Dalam domain AIED, Memarian \& Doleck [5] menekankan bahwa implementasi HITL yang bermakna memerlukan pencatatan data interaksi yang komprehensif—tanpa data yang memadai, pola perilaku pengguna tidak dapat diidentifikasi dan peran manusia dalam loop tidak dapat dianalisis secara sistematis. Pencatatan yang terstruktur juga merupakan prasyarat untuk reproduktibilitas analisis: peneliti lain harus dapat mereplikasi proses analisis berdasarkan data \textit{log} yang dikumpulkan [5].

Dalam proyek \textit{Escape the Sketchbook}, \textit{data logging} dirancang untuk menangkap konteks lengkap setiap keputusan siswa terhadap output AI. Setiap entri \textit{log} mencatat beberapa \textit{field} yang masing-masing memiliki rasional tersendiri. Pertama, \texttt{sketch\_id} mengidentifikasi sketsa yang diklasifikasikan, memungkinkan analisis berdasarkan tingkat kesulitan atau ambiguitas sketsa tertentu. Kedua, \texttt{top3\_predictions} dan \texttt{top3\_confidences} mencatat tiga prediksi teratas beserta nilai \textit{confidence}-nya, memberikan konteks tentang informasi apa yang tersedia bagi siswa saat membuat keputusan. Ketiga, \texttt{user\_decision} mencatat apakah siswa setuju atau menolak prediksi utama AI, yang merupakan variabel utama dalam analisis pola keputusan. Keempat, \texttt{response\_time} mencatat durasi antara saat \textit{confidence score} ditampilkan dan saat siswa mengambil keputusan, memberikan indikasi tentang seberapa banyak waktu yang dihabiskan siswa untuk mengevaluasi output AI. Kelima, \texttt{session\_id} dan \texttt{timestamp} memungkinkan pengelompokan \textit{log} berdasarkan sesi dan analisis temporal.

Arsitektur \textit{logging} dalam proyek ini menggunakan \textit{REST API} sebagai antarmuka antara \textit{client} (\textit{browser} siswa) dan server. Ketika siswa membuat keputusan, \textit{client} mengirimkan data \textit{log} ke server melalui permintaan HTTP POST, dan server menyimpan data tersebut ke \textit{database} SQLite. Pemilihan SQLite sebagai \textit{database} didasarkan pada kesederhanaan dan kecukupan untuk skala data yang diharapkan—SQLite tidak memerlukan konfigurasi \textit{server database} terpisah dan sudah memadai untuk penyimpanan data \textit{log} dalam konteks proyek akhir. Penting untuk ditekankan bahwa server dalam arsitektur ini hanya berperan sebagai penyimpan \textit{log}—seluruh inferensi model berjalan di \textit{client-side} menggunakan TensorFlow.js.

% ============================================================
\subsection{K-Means Clustering}
\label{subsec:kmeans}
% ============================================================

\textit{K-Means clustering} merupakan algoritma \textit{unsupervised learning} yang bertujuan mengelompokkan data ke dalam $k$ kelompok (\textit{cluster}) berdasarkan kemiripan fitur, di mana setiap data diatribusikan ke kelompok dengan \textit{centroid} terdekat. Algoritma ini bekerja secara iteratif: (1) inisialisasi $k$ \textit{centroid} secara acak, (2) atribusikan setiap data ke \textit{centroid} terdekat, (3) perbarui posisi \textit{centroid} berdasarkan rata-rata seluruh data dalam kelompok, dan (4) ulangi langkah 2--3 hingga konvergensi. Meskipun sederhana, K-Means telah terbukti efektif untuk analisis pola dalam berbagai domain, termasuk pendidikan [10][11].

Pansri \textit{et al.} [10] menerapkan K-Means untuk mengidentifikasi pola perilaku belajar siswa berdasarkan data interaksi pada \textit{learning management system}. Mereka menggunakan fitur-fitur seperti frekuensi akses, durasi sesi, dan pola penyelesaian tugas untuk mengelompokkan siswa ke dalam beberapa profil perilaku. Temuan mereka menunjukkan bahwa K-Means mampu menghasilkan kelompok yang dapat diinterpretasikan secara bermakna—misalnya, kelompok siswa yang aktif dan konsisten versus kelompok siswa yang sporadis—meskipun hasil pengelompokan bersifat deskriptif dan tidak mengimplikasikan kausalitas [10]. Alzahrani \textit{et al.} [11] menggunakan pendekatan serupa untuk menganalisis \textit{engagement} mingguan siswa dan menemukan bahwa K-Means dapat mengidentifikasi pola \textit{engagement} yang berbeda antar kelompok siswa, yang dapat digunakan sebagai dasar intervensi pedagogis.

Dalam konteks proyek \textit{Escape the Sketchbook}, K-Means diterapkan untuk menganalisis pola keputusan siswa berdasarkan data \textit{log} interaksi. Fitur-fitur yang digunakan untuk pengelompokan mencakup: rasio persetujuan terhadap prediksi AI (seberapa sering siswa menyetujui prediksi utama), rata-rata \textit{confidence gap} pada keputusan yang disetujui versus ditolak, dan rata-rata \textit{response time}. Berdasarkan konfigurasi fitur ini, hasil pengelompokan K-Means diinterpretasikan ke dalam tiga profil keputusan: (1) \textit{Trust Acceptor}—siswa yang cenderung menyetujui prediksi AI tanpa mempertimbangkan alternatif, ditandai dengan rasio persetujuan tinggi dan \textit{response time} rendah; (2) \textit{Critical Evaluator}—siswa yang mengevaluasi \textit{confidence score} sebelum mengambil keputusan, ditandai dengan variasi rasio persetujuan yang bergantung pada \textit{confidence gap} dan \textit{response time} yang lebih tinggi; serta (3) \textit{Uncertain Explorer}—siswa yang menunjukkan pola keputusan tidak konsisten, ditandai dengan variasi tinggi pada seluruh metrik.

Penting untuk ditekankan bahwa hasil K-Means dalam proyek ini bersifat \textbf{interpretif, bukan diagnostik}. Pengelompokan K-Means mengidentifikasi pola-pola empiris dalam data keputusan, namun tidak memberikan diagnosis tentang kemampuan kognitif atau tingkat literasi AI siswa. Profil \textit{Trust Acceptor}, \textit{Critical Evaluator}, dan \textit{Uncertain Explorer} merupakan label interpretif yang didasarkan pada pengamatan pola data, bukan klasifikasi klinis atau penilaian kompetensi. Interpretasi ini bertujuan untuk memperlihatkan variasi pola keputusan di antara siswa, yang dapat menjadi bahan refleksi bagi guru—bukan untuk mengkategorikan siswa ke dalam hierarki kemampuan. Selain itu, jumlah kelompok ($k$) ditentukan melalui metode \textit{elbow method} dan validasi menggunakan \textit{silhouette score}, memastikan bahwa pengelompokan yang dihasilkan memiliki kohesi internal yang memadai.

% ============================================================
\section{PENELITIAN TERKAIT}
\label{sec:penelitian-terkait}
% ============================================================

Bagian ini membahas penelitian-penelitian terkait yang menjadi referensi dalam perancangan komponen teknis proyek \textit{Escape the Sketchbook}. Setiap penelitian dianalisis dari segi relevansi dan kontribusinya terhadap aspek-aspek spesifik yang dikembangkan dalam proyek ini.

% ============================================================
\subsection{Deep Learning for Free-Hand Sketch Recognition}
\label{subsec:penelitian-sketch}
% ============================================================

Xu \textit{et al.} [6] menyajikan survei komprehensif mengenai penerapan \textit{deep learning} untuk pengenalan sketsa tangan bebas. Kajian mereka mencakup evolusi arsitektur dari metode berbasis \textit{hand-crafted features} (seperti HOG dan SIFT) hingga pendekatan \textit{end-to-end deep learning} menggunakan CNN. Salah satu temuan utama adalah bahwa arsitektur CNN standar seperti VGG dan ResNet, meskipun sangat efektif untuk klasifikasi gambar fotografis, tidak selalu menjadi pilihan optimal untuk sketsa karena karakteristik data yang berbeda. Sketsa bersifat \textit{sparse} dan biner (garis pada latar belakang putih), sehingga representasi fitur yang dipelajari dari gambar fotografis berwarna tidak sepenuhnya transferabel. Xu \textit{et al.} juga mengevaluasi berbagai teknik \textit{data augmentation} khusus untuk sketsa—seperti \textit{random scaling}, \textit{rotation}, dan \textit{stroke perturbation}—yang terbukti meningkatkan performa klasifikasi secara signifikan [6].

Relevansi penelitian ini terhadap proyek \textit{Escape the Sketchbook} terletak pada pemahaman bahwa klasifikasi sketsa memerlukan penanganan yang berbeda dari klasifikasi foto. Temuan Xu \textit{et al.} mengenai sifat \textit{sparse} dan ambiguitas sketsa mendukung keputusan desain untuk menampilkan \textit{Top-3 confidence score}—karena ambiguitas inheren sketsa, model seringkali tidak dapat memberikan prediksi dengan \textit{confidence} tinggi pada satu kelas, sehingga menampilkan beberapa alternatif merupakan representasi yang lebih jujur terhadap kondisi model. Selain itu, teknik \textit{data augmentation} yang mereka identifikasi menjadi referensi untuk pra-pemrosesan data latih pada model MobileNet yang digunakan dalam proyek ini. Namun, Xu \textit{et al.} tidak membahas aspek pedagogis atau penggunaan klasifikasi sketsa sebagai sarana literasi AI—yang merupakan kontribusi utama proyek ini.

% ============================================================
\subsection{Front-End Deep Learning Web Applications}
\label{subsec:penelitian-frontend-dl}
% ============================================================

Goh, Ho \& Abas [7] melakukan studi mendalam mengenai penerapan \textit{deep learning} pada \textit{front-end web application}, dengan fokus pada arsitektur, kerangka kerja, dan tantangan \textit{deployment}. Mereka mengidentifikasi bahwa pendekatan \textit{front-end deployment}—di mana model berjalan langsung di \textit{browser} menggunakan \textit{library} seperti TensorFlow.js atau ONNX.js—memiliki keunggulan dalam hal latensi (tidak ada \textit{round-trip} ke server), privasi data (data tidak keluar dari perangkat pengguna), dan aksesibilitas (tidak memerlukan infrastruktur server). Namun, mereka juga mencatat tantangan signifikan: keterbatasan memori \textit{browser}, variasi dukungan WebGL antar \textit{browser}, dan \textit{trade-off} antara ukuran model dan akurasi yang harus dikelola dengan cermat [7].

Goh \textit{et al.} secara spesifik mengevaluasi performa beberapa arsitektur model pada lingkungan \textit{browser}, dan menemukan bahwa MobileNetV2 dengan \textit{width multiplier} $\alpha = 0.5$ atau $\alpha = 0.75$ memberikan keseimbangan terbaik antara akurasi klasifikasi dan kebutuhan sumber daya—waktu inferensi berkisar antara 30--100 ms pada perangkat menengah, dengan ukuran model di bawah 5 MB [7]. Temuan ini sangat relevan dengan keputusan arsitektural dalam proyek \textit{Escape the Sketchbook} untuk menggunakan MobileNet sebagai model klasifikasi yang dijalankan di \textit{client-side}. Perbedaan mendasar antara kajian Goh \textit{et al.} dan proyek ini adalah bahwa Goh \textit{et al.} berfokus pada aspek teknis \textit{deployment}, sementara proyek ini mengintegrasikan aspek \textit{deployment front-end} dengan mekanisme HITL, \textit{data logging}, dan analisis \textit{K-Means clustering} dalam satu sistem yang koheren.

% ============================================================
\subsection{Confidence Visualization}
\label{subsec:penelitian-confidence-vis}
% ============================================================

Karran, Demazure \& Hudelot [9] mengkaji desain visualisasi \textit{confidence} untuk membantu pengguna menilai dan menimbang output sistem AI. Melalui analisis literatur dan eksperimen desain, mereka mengidentifikasi bahwa efektivitas visualisasi \textit{confidence} bergantung pada tiga faktor: (1) \textit{granularity}—apakah visualisasi menampilkan satu nilai atau distribusi lengkap, (2) \textit{comparability}—apakah visualisasi memungkinkan perbandingan antar alternatif, dan (3) \textit{contextual framing}—apakah visualisasi menyertakan konteks yang membantu pengguna menginterpretasikan nilai \textit{confidence}. Mereka menemukan bahwa visualisasi yang menampilkan distribusi \textit{confidence} di beberapa alternatif—misalnya \textit{bar chart} horizontal untuk \textit{Top-N prediksi}—lebih efektif dalam memfasilitasi evaluasi kritis dibandingkan visualisasi yang hanya menampilkan satu angka [9].

Relevansi penelitian Karran \textit{et al.} terhadap proyek ini terletak pada dukungan empiris untuk keputusan menampilkan \textit{Top-3 confidence score} alih-alih hanya label prediksi utama. Prinsip \textit{comparability} yang mereka identifikasi menegaskan bahwa siswa perlu dapat membandingkan beberapa alternatif untuk mengevaluasi output AI secara bermakna—yang merupakan esensi dari mekanisme HITL dalam sistem ini. Selain itu, temuan mereka mengenai \textit{contextual framing} mendukung keputusan untuk menyertakan kategori subset (\textit{Solid}/\textit{Danger}/\textit{Decorative}) sebagai konteks tambahan yang membantu siswa menginterpretasikan prediksi model. Namun, Karran \textit{et al.} tidak membahas penggunaan visualisasi \textit{confidence} dalam konteks literasi AI untuk siswa SMP, dan juga tidak mengintegrasikan visualisasi \textit{confidence} dengan mekanisme \textit{logging} dan analisis \textit{clustering}—yang merupakan kontribusi integratif proyek ini.

% ============================================================
\subsection{K-Means untuk Pola Belajar}
\label{subsec:penelitian-kmeans}
% ============================================================

Pansri \textit{et al.} [10] menerapkan algoritma K-Means untuk menganalisis pola perilaku belajar siswa berdasarkan data interaksi pada \textit{learning management system} (LMS). Fitur-fitur yang digunakan mencakup frekuensi akses materi, durasi sesi belajar, pola penyelesaian tugas, dan skor kuis. Hasil pengelompokan mengidentifikasi empat profil perilaku: \textit{active learners} (siswa yang aktif dan konsisten), \textit{sporadic learners} (siswa yang mengakses secara tidak teratur), \textit{struggling learners} (siswa yang aktif namun berkinerja rendah), dan \textit{disengaged learners} (siswa yang jarang mengakses dan berkinerja rendah). Pansri \textit{et al.} menekankan bahwa profil ini bersifat deskriptif—menggambarkan pola observasi—dan tidak mengimplikasikan diagnosis atau prediksi tentang kemampuan siswa [10].

Alzahrani \textit{et al.} [11] menggunakan pendekatan serupa untuk menganalisis \textit{engagement} mingguan siswa. Mereka mengumpulkan data partisipasi, waktu yang dihabiskan, dan interaksi dengan konten selama satu semester, kemudian mengelompokkan siswa menggunakan K-Means. Temuan mereka menunjukkan bahwa pola \textit{engagement} dapat berubah sepanjang semester—siswa yang awalnya aktif dapat menjadi pasif, dan sebaliknya—yang menunjukkan pentingnya analisis temporal dalam interpretasi hasil \textit{clustering}. Alzahrani \textit{et al.} juga mengevaluasi metode penentuan jumlah kelompok ($k$) dan menemukan bahwa kombinasi \textit{elbow method} dengan \textit{silhouette score} menghasilkan pengelompokan yang paling stabil dan dapat diinterpretasikan [11].

Relevansi kedua penelitian ini terhadap proyek \textit{Escape the Sketchbook} terletak pada metodologi penerapan K-Means untuk analisis pola perilaku dalam konteks pendidikan. Pansri \textit{et al.} dan Alzahrani \textit{et al.} menunjukkan bahwa K-Means mampu menghasilkan profil perilaku yang dapat diinterpretasikan dari data interaksi—yang menjadi dasar bagi penerapan K-Means pada data \textit{log} keputusan siswa dalam proyek ini. Perbedaan fundamental adalah bahwa kedua penelitian tersebut menganalisis perilaku belajar umum (akses LMS, penyelesaian tugas), sementara proyek ini menganalisis pola keputusan spesifik terhadap output AI—yaitu bagaimana siswa merespons \textit{confidence score} dan apakah mereka mengevaluasi alternatif sebelum mengambil keputusan. Domain analisis yang berbeda ini menghasilkan profil yang secara konseptual berbeda: \textit{Trust Acceptor}, \textit{Critical Evaluator}, dan \textit{Uncertain Explorer} merupakan profil keputusan terhadap AI, bukan profil gaya belajar umum.

% ============================================================
\subsection{State of The Art}
\label{subsec:state-of-the-art}
% ============================================================

Berdasarkan kajian terhadap penelitian-penelitian terkait, dapat diidentifikasi posisi \textit{state of the art} dan celah penelitian yang diisi oleh proyek \textit{Escape the Sketchbook}. Tabel~\ref{tab:sota} merangkum cakupan masing-masing penelitian terkait.

\begin{table}[H]
\centering
\caption{Rangkuman cakupan penelitian terkait}
\label{tab:sota}
\small
% Catatan: tambahkan \usepackage{rotating} pada preamble dokumen utama
\begin{tabular}{|l|c|c|c|c|c|c|}
\hline
\textbf{Penelitian} & \rotatebox{90}{\textbf{Klasifikasi Sketsa}} & \rotatebox{90}{\textbf{Client-Side Inference}} & \rotatebox{90}{\textbf{Confidence Vis.}} & \rotatebox{90}{\textbf{HITL Logging}} & \rotatebox{90}{\textbf{K-Means Analisis}} & \rotatebox{90}{\textbf{Integrasi Sistem}} \\
\hline
Xu \textit{et al.} [6] & \checkmark & -- & -- & -- & -- & -- \\
\hline
Goh \textit{et al.} [7] & -- & \checkmark & -- & -- & -- & -- \\
\hline
Karran \textit{et al.} [9] & -- & -- & \checkmark & -- & -- & -- \\
\hline
Pansri \textit{et al.} [10] & -- & -- & -- & -- & \checkmark & -- \\
\hline
Alzahrani \textit{et al.} [11] & -- & -- & -- & -- & \checkmark & -- \\
\hline
\textbf{Proyek ini} & \checkmark & \checkmark & \checkmark & \checkmark & \checkmark & \checkmark \\
\hline
\end{tabular}
\end{table}

Dari tabel di atas, terlihat bahwa setiap penelitian terkait memberikan kontribusi pada satu aspek spesifik: Xu \textit{et al.} [6] pada klasifikasi sketsa, Goh \textit{et al.} [7] pada \textit{deployment front-end}, Karran \textit{et al.} [9] pada visualisasi \textit{confidence}, serta Pansri \textit{et al.} [10] dan Alzahrani \textit{et al.} [11] pada analisis \textit{K-Means} untuk pola perilaku. Namun, tidak ada satu pun dari penelitian tersebut yang mengintegrasikan seluruh komponen—klasifikasi sketsa, inferensi \textit{client-side}, visualisasi \textit{confidence}, mekanisme HITL dengan \textit{logging} terstruktur, dan analisis \textit{K-Means clustering}—dalam satu sistem yang koheren.

Celah penelitian yang diisi oleh proyek ini memiliki karakteristik sebagai berikut. Pertama, integrasi antara klasifikasi sketsa dan visualisasi \textit{confidence}: sementara Xu \textit{et al.} [6] berfokus pada akurasi klasifikasi dan Karran \textit{et al.} [9] berfokus pada desain visualisasi, proyek ini menghubungkan keduanya dengan menjadikan distribusi \textit{confidence} dari klasifikasi sketsa sebagai objek evaluasi oleh siswa. Kedua, integrasi antara HITL dan \textit{data logging}: sementara Memarian \& Doleck [5] mendiskusikan peran HITL dalam AIED secara konseptual, proyek ini mengimplementasikan HITL sebagai mekanisme pencatatan keputusan yang menghasilkan data terstruktur untuk analisis. Ketiga, integrasi antara \textit{logging} dan \textit{K-Means clustering}: sementara Pansri \textit{et al.} [10] dan Alzahrani \textit{et al.} [11] mengaplikasikan K-Means pada data LMS, proyek ini mengaplikasikan K-Means pada data keputusan siswa terhadap output AI—sebuah domain analisis yang secara konseptual berbeda.

Kontribusi teknis Muhammad Dias Al-Izzat dalam proyek ini secara spesifik mencakup: (1) perancangan dan pelatihan model CNN berbasis MobileNet untuk klasifikasi sketsa dengan \textit{transfer learning}, (2) implementasi inferensi \textit{client-side} menggunakan TensorFlow.js di \textit{browser}, (3) perancangan skema \textit{data logging} terstruktur melalui \textit{REST API} dan penyimpanan ke SQLite, (4) implementasi analisis \textit{K-Means clustering} untuk mengidentifikasi pola keputusan siswa, serta (5) pengembangan \textit{dashboard} untuk visualisasi hasil analisis. Kelima komponen ini diintegrasikan dalam satu arsitektur sistem di mana inferensi berjalan di \textit{client-side}, keputusan dicatat melalui \textit{REST API}, dan analisis dilakukan secara \textit{post-hoc} menggunakan data \textit{log}—menciptakan alur data yang koheren dari interaksi siswa hingga interpretasi pola keputusan.
